\section{Exact Proof Boundary}
\label{sec:exact-proof-boundary}

The verified results above do not form an unconditional operational
theorem. This section states the boundary as part of the theorem.

\begin{itemize}
  \item \textbf{Square bound.} \texttt{SpecSieveSequence.next} and the
    next-head theorem require $p^+\lt h^2$. Bertrand's postulate
    supplies a prime between $h$ and $2h$; for prime $h\ge2$, this
    implies the required square bound (with $h=2$ checked directly).
    Ramanujan gives a direct proof of the required postulate
    \cite{ramanujan1919}, but Bertrand's postulate is an external
    mathematical dependency here, not a Stainless theorem in this
    development.

  \item \textbf{Count-to-period bridge.} Section~5.1 verifies that
    filtering the complete expanded current-stage window leaves exactly
    $T(h-1)$ values. Section~6.3 separately establishes the semantic
    gap-prefix agreement, while Section~6.4 reconstructs the next stage
    when $\ell'_{T'}=h'+M'$ is supplied. The derivation of that next
    canonical-period equation from the same-head count is a distinct
    open composition theorem.

  \item \textbf{Direct cycle-to-cycle construction.} Repetition
    preserves the represented values; filtering the base and repeated
    views gives equal survivor lists and gaps; the first survivor is the
    next head; and the semantic merged-gap prefix agrees with the next
    specification. The open composition theorem is the equality between
    the repeated cycle's filtered survivor-gap list and that semantic
    merged-gap prefix, followed by packaging those gaps into a new
    integral cycle.

  \item \textbf{No short-window gap-persistence theorem.} Section~5.2
    proves that exactly $h-2$ lifts of one real linear 2-gap keep both
    endpoints over a complete lift block. It does not imply that one of
    those lifts lies in a shorter interval such as $[h,h^2)$, nor does
    it yet aggregate the cyclic wrap gap into a total next-stage
    recurrence. This article makes no claim about infinitely many twin
    primes or the survival of 2-gaps in every local window.

  \item \textbf{No efficiency theorem.} The period grows from $T$ to
    $T(h-1)$. The finite cycle is analytically useful, but materializing
    it need not outperform a conventional segmented sieve. No time or
    space complexity advantage is claimed.
\end{itemize}

These qualifications are part of the result: they separate the verified
finite-stage theorem from adjacent mathematical questions.
